\clearpage
\chapter*{2025 Septiembre}
\addcontentsline{toc}{chapter}{2025 Septiembre}

\begin{enunciado}
    \hypertarget{2025 S 1}{}
    \textbf{1. (2 puntos)} Demostrar que un espacio topológico es compacto si y solo si se verifica que en él toda base de filtro \(\{B_\lambda\}_{\lambda \in L}\) tiene un punto de aglomeración.
\end{enunciado}

\begin{enunciado}
    \hypertarget{2025 S 2}{}
    \textbf{2. (4 puntos)} En el conjunto \(N\) de los números naturales (se incluye el 0), definimos la topología \(T = \{A \subset N \mid \text{si } 2n + 1 \notin A \text{ entonces } 2n, 2n + 2 \notin A\}\).
    \begin{enumerate}[label=\alph*)]
        \item Encontrar una base de \(T\).
        \item Estudiar si \((N,T)\) es \(T_2\).
    \end{enumerate}
\end{enunciado}

\begin{enunciado}
    \hypertarget{2025 S 3}{}
    \textbf{3. (4 puntos)}
    \begin{enumerate}[label=\alph*)]
        \item Sean \((X,T)\) un espacio topológico, \(A \subset X\) y \(r : X \to A\) una aplicación continua de \((X,T)\) en un espacio topológico \((A,T_A)\) tal que \(r/A = 1_A\) (la identidad en \(A\)). Estudiar qué relación hay entre \(T_A\) y la topología final en \(A\) para la aplicación \(r\), de la topología \(T\).
        \item En el conjunto \(\N\) de los números naturales (se incluye el 0), se considera la topología
        \[
            T = \{M \subset \N \mid \text{si } n \in M \text{ entonces todo divisor de } n \text{ pertenece a } M\}.
        \]
        Estudiar si \((X,T)\) es compacto.
    \end{enumerate}
\end{enunciado}